\section*{\textbf{ABSTRACT}}
\addcontentsline{toc}{section}{ABSTRACT}

Birds are a key part of the ecosystems worldwide. They help keep the balance in the ecosystem and act as powerful bioindicators of ecosystem health. This is why monitoring bird populations on a large-scale is important. However, traditional methods to keep track of species rely on manual observations and expert knowledge. This makes large-scale monitoring impractical as it is very time consuming and difficult to scale. This project proposes BirdSight, a web-based citizen science platform built using Java Spring Boot for the backend and Next.js for the frontend. The platform integrates a machine learning model for species recognition, using deep learning techniques to identify bird species from input images. Multiple convolutional neural network architectures, including VGG16, ResNet50, InceptionV3, MobileNetV2, and DenseNet121, were trained and evaluated. The best-performing and most suitable model for deployment was then further optimized and integrated into the platform. The platform aims to enhance large-scale bird monitoring, promote citizen science, and encourage public engagement with biodiversity around them through a scalable digital platform. The platform allows users to submit observations, visualize sightings on interactive maps, become a part of a community of like-minded people, and contribute to biodiversity data collection.

\vspace{1em}
\noindent\textbf{Keywords:} Bird Species Classification, Citizen Science, Convolutional Neural Networks, Deep Learning, Image Classification, Biodiversity Monitoring, Web Platform, Comparative Analysis

\clearpage